\documentclass[12pt,a4paper]{article}
\usepackage[utf8]{inputenc}
\usepackage[T2A]{fontenc}
\usepackage[russian]{babel}
\usepackage{amsmath,amssymb}
\usepackage{geometry}
\geometry{left=2cm,right=2cm,top=2cm,bottom=2cm}

\author{Иван Насекин БКНАД252}
\title{Лабораторная работа C++. \\ Исследование разных типов данных. Сравнение производительности битовых и арифметических операций.}
\date{22 января 2026 г.}

\begin{document}

    \maketitle

    \section*{Task 01. Битовые операции для чисел}

    Результаты:

    a \& b = 42

    a | b = -17

    a \^{} b = -59

    \~{}a = -43

    a << 2 = 168, a >> 2 = 10

    b << 2 = -68, b >> 2 = -5

    Двоичное представление:

    42: 00000000000000000000000000101010

    -17: 11111111111111111111111111101111

    \subsection*{Чему равно -1 в 32-битном представлении?}

    -1 это 11111111111111111111111111111111 (все единицы).

    \subsection*{Как устроен дополнительный код?}

    Чтобы получить отрицательное число в дополнительном коде, берем положительное число, инвертируем все биты и прибавляем 1. Например, для -17: сначала 17 = ...010001, потом инвертируем = ...101110, и +1 = ...101111.
    Дополнительный код нужен потому что с ним процессору проще делать арифметику — вычитание превращается в сложение и получается что процессору не нужна отдельная схема для вычитания.

    \subsection*{Почему сдвиг отрицательных чисел дает отрицательный результат?}

    Когда сдвигаем отрицательное число вправо, компьютер делает арифметический сдвиг — слева добавляются единицы вместо нулей, чтобы сохранить знак. Поэтому число остается отрицательным.

    \section*{Вывод}

    Разобрались, как числа хранятся на уровне битов. Отрицательные числа используют дополнительный код, чтобы процессор мог складывать их так же, как обычные. Битовые операции работают напрямую с двоичным представлением.

    \section*{Task 02. Переполнение целочисленных типов}

    UINT32\_MAX + 2 = 1 \\
    INT32\_MAX + 2 = -2147483647

    \subsection*{UINT32\_MAX + 2}

    UINT32\_MAX = 4294967295, в двоичном виде все 32 бита равны единице. Когда прибавляем 2, происходит переполнение:

    \begin{verbatim}
  11111111111111111111111111111111  (4294967295)
+ 00000000000000000000000000000010  (+2)
----------------------------------
 100000000000000000000000000000001
    \end{verbatim}

    Поскольку uint32\_t хранит только 32 бита, старший 33-й бит отбрасывается, остается 00000000000000000000000000000001, что равно 1.

    \subsection*{INT32\_MAX + 2}

    INT32\_MAX = 2147483647, в двоичном виде старший бит 0, остальные 31 бит равны 1. При добавлении 2:

    \begin{verbatim}
  01111111111111111111111111111111  (2147483647)
+ 00000000000000000000000000000010  (+2)
----------------------------------
  10000000000000000000000000000001
    \end{verbatim}

    Старший бит стал 1, что означает отрицательное число в дополнительном коде. Это -2147483647.

    \subsection*{Почему переполнение беззнаковых типов безопасно, а знаковых — нет?}

    С беззнаковыми типами все просто — когда число выходит за максимум, оно просто заворачивается обратно к нулю. Стандарт C++ это четко прописывает, поэтому результат всегда предсказуемый.

    А вот со знаковыми типами проблемы — переполнение это ошибка undefined behavior. То есть стандарт вообще не говорит что должно случиться. Компилятор может решить что переполнения не будет никогда и сделать какие-то свои оптимизации, из-за чего программа начнет вести себя непредсказуемо.

    \subsection*{Что говорит стандарт C++ о переполнении знаковых целых?}

    Стандарт говорит что это undefined behavior. Короче, никаких гарантий нет — может упасть, может выдать какую-то ерунду, а может и работать нормально. Зависит от компилятора и настроек.

    \subsection*{Как избежать переполнения в реальных программах?}

    \begin{itemize}
        \item Проверки через if/else
        \item Если число точно положительное, можно взять беззнаковый тип
        \item Использовать типы с большим размером
        \item Включить проверки компилятора типа
    \end{itemize}

    \section*{Вывод}

    Изучили переполнение целых чисел. Беззнаковые типы переполняются предсказуемо по модулю, знаковые приводят к ошибке переполнения. Нужно проверять входные данные и по возможности использовать подходящие типы для предотвращения переполнений.

    \section*{Task 03. Особенности вещественных чисел}

    \subsection*{Представление числа 10.25}

    10.25 в десятичной = 1010.01 в двоичной = $1.01001 \times 2^3$\\
    Знак: 0 (положительное)\\
    Экспонента: $3 + 127 = 130 = 10000010_2$\\
    Мантисса: 01001000000000000000000\\

    \subsection*{Представление числа 7.1}

    7.1 в двоичной это бесконечная периодическая дробь.

    \subsection*{Сумма 10.25 + 7.1}

    Результат с 12 знаками: 17.3500003815

    Должно было получится 17.35, но получили чуть больше. Причина в 7.1 нельзя точно записать в двоичной системе, поэтому накапливается ошибка округления.

    \subsection*{Нарушение ассоциативности}

    float a = 16'777'216.0f;\\
    float b = 1.0f;\\
    float c = 1.0f;\\
    (a + b) + c = 16777216 \\
    a + (b + c) = 16777218

    Различие в результатах происходит из-за ограниченной точности типа float и необходимости округления при арифметических операциях.

    \subsection*{Сравнение на равенство вещественных чисел}

    Основная причина заключается в том, что большинство десятичных дробей невозможно представить в двоичной системе счисления точно.

    Привести может к непредсказумемым ошибкам, к трудноуловимым багам.

    \section*{Вывод}
    Обнаружили, что многие десятичные дроби невозможно точно представить в двоичной системе, что приводит к ошибкам округления. Ассоциативность арифметических операций не гарантируется из-за конечной точности. Прямое сравнение вещественных чисел на равенство является плохой практикой программирования и может приводить к логическим ошибкам. Для надежной работы с вещественными числами необходимо использовать сравнение с допустимой погрешностью и понимать ограничения точности используемого типа данных.

    \section*{Task 04. Битовые vs арифмитические операции}

    Битовые операции быстрее арифметических в случае деления знаковых чисел.\\
    Деление и модуль — сложные операции процессора, особенно для знаковых чисел. Битовые операции являются простейшими инструкциями, выполняются за 1-2 такта.\\
    Для беззнакового деления и умножения разница минимальна.

    \section*{Вывод}
    Провели сравнение производительности битовых и арифметических операций. Битовые операции показали значительное преимущество при делении знаковых чисел, что объясняется сложностью инструкции деления на уровне процессора. Для умножения и деления на степени двойки компилятор способен автоматически заменять арифметические операции на битовые сдвиги, особенно для беззнаковых типов.

    \section*{Task 05. Влияние типа данных на скорость вычислений}

    \subsection*{Как отличаются результаты с volatile и без?}
    Результаты с volatile и без почти одинаковые. volatile запрещает оптимизации, но здесь оптимизаций и так нет или они минимальны.

    \subsection*{Зависит ли время выполнения от размера типа? Почему?}
    Да, зависит, int в 2 раза быстрее char/short. Процессор заточен под 32-битные операции, а с 8/16-битными типами приходится делать дополнительные действия например обрезка битов, расширение знака.

    \subsection*{Как ведут себя float и double по сравнению с целыми?}
    Примерно в 3.5 раза медленнее. Вещественная арифметика сложнее тк нужно работать с мантиссой, экспонентой, округлением. Это делает отдельный блок FPU, а не основной ALU

    \subsection*{Влияет ли знаковость у целых типов на производительность?}
    Практически нет. Разница между signed и unsigned в пределах погрешности.\\
    Для сложения знаковость не важна - операция одинаковая на уровне процессора.

    \section*{Вывод}
    int быстрее всех, char/short медленнее в 2 раза, float/double - в 3.5 раза. Процессор лучше работает с 32-битными целыми числами, для которых он и оптимизирован.

\end{document}